\chapter{TƯ DUY PHẢN BIỆN ĐỂ ĐÁNH GIÁ ĐỘ XÁC THỰC CỦA DỮ LIỆU}
\label{ch:M5L4LessonNotes}

\section*{Lesson Notes}

MODULE 5 \textbar{} LESSON 4

Thời gian đọc: 15 phút

Kiến thức nền: Toàn bộ nội dung Financial Data, nội dung Financial Markets Module 7

Từ khóa: Năng lực thông tin (Information Literacy), Tin giả (Fake News)

Trong bài học này, chúng ta xem xét các đặc điểm của tư duy phản biện (critical thinking). Chúng ta sẽ tìm hiểu một bài tổng quan tài liệu có hệ thống, trong đó đề cập đến việc các kỹ năng tư duy phản biện được sử dụng như thế nào để xác định một tin tức có phải là "tin giả" hay không.

\section{Tư duy phản biện}

Tư duy phản biện không chỉ là điều chúng ta bàn một lần rồi thôi. Như Ismail viết: "Tư duy phản biện bao gồm tư duy có chủ đích, hướng đến mục tiêu trong quá trình đưa ra quyết định dựa trên bằng chứng trong quá trình giải quyết vấn đề khoa học" (1). Trong suốt chương trình này, bạn sẽ được kỳ vọng vận dụng các kỹ năng tư duy phản biện cho nhiều mục đích khác nhau, nhưng ở đây, chúng ta sẽ tập trung vào việc dùng tư duy phản biện để xác định một nguồn tin có đáng tin cậy hay không.

Bài báo của Ismail đề cập tầm quan trọng của các kỹ năng tư duy phản biện đối với việc được tuyển dụng và làm việc hiệu quả sau khi đã có việc làm. Bài báo nêu bật một nghiên cứu được thực hiện trên những nhân viên có năm năm kinh nghiệm trong công việc kỹ thuật, đặc biệt là kỹ thuật cơ khí. Khoảng thời gian năm năm được xem là điểm khởi đầu để hình thành sự trưởng thành trong sự nghiệp của một người.

Theo nghiên cứu này, bảng câu hỏi sau đã được gửi cho các kỹ sư cơ khí có 5 năm kinh nghiệm làm việc. Các câu hỏi liên quan đến:

\begin{itemize}
\item
  Xác định vấn đề trong một hệ thống phức tạp
\item
  Cải thiện kỹ năng giải thích, phân tích và đánh giá
\item
  Tìm kiếm các giải pháp thay thế
\item
  Tìm kiếm các giải pháp sáng tạo
\item
  Biện minh cho các quyết định dựa trên bằng chứng vững chắc
\item
  Khả năng kiên trì hoàn thành các trách nhiệm được giao
\item
  Khả năng hiểu và thích nghi với văn hóa của cộng đồng
\end{itemize}

Các kỹ năng tư duy phản biện rất được ưa chuộng trong môi trường làm việc vì chúng cần thiết để làm việc hiệu quả trong nhóm, phân tích dữ liệu, xây dựng lập luận, ra quyết định và giải quyết vấn đề. Do đó, điều quan trọng là bạn nên củng cố kỹ năng này hết mức có thể trong suốt chương trình.

Kỹ năng tư duy phản biện của bạn sẽ được đánh giá theo nhiều cách khác nhau, từ các bài đăng trên diễn đàn, các dự án làm việc nhóm cho đến dự án tốt nghiệp (capstone). Một kỹ năng tư duy phản biện mà chúng ta tập trung ở đây là năng lực thông tin (information literacy), vốn rất cần thiết khi bạn tìm kiếm các nguồn hợp lệ và đáng tin cậy để tìm dữ liệu cần dùng, để củng cố và phát triển ý tưởng của mình, v.v. Những kỹ năng này vô cùng quan trọng trong các công việc kỹ thuật tài chính cũng như các ngành nghề liên quan đến STEM khác. Ở phần sau của bài học này, chúng ta sẽ vận dụng kỹ năng tư duy phản biện để xác định một tin tức là tin giả hay đáng tin cậy.

\begin{enumerate}
\def\labelenumi{\arabic{enumi}.}
\setcounter{enumi}{1}
\item
  Năng lực thông tin
\end{enumerate}

Hiệp hội Thư viện Hoa Kỳ (American Library Association) định nghĩa năng lực thông tin là "một tập hợp các khả năng đòi hỏi mỗi cá nhân phải nhận biết khi nào cần thông tin và có khả năng tìm kiếm, đánh giá và sử dụng hiệu quả thông tin cần thiết." Như vậy, năng lực thông tin nghĩa là trước hết bạn biết cách nghiên cứu và tìm các nguồn mình cần, đồng thời biết cách đánh giá xem một nguồn có phải là nguồn thông tin đáng tin cậy và xác thực hay không. Tuy nghe có vẻ đơn giản, nhưng đây đã trở thành một nhiệm vụ ngày càng khó khăn khi tin giả, thông tin do cộng đồng đóng góp (crowdsourced) và các tác giả ẩn danh trên internet xuất hiện.

Vì vậy, việc đầu tiên chúng tôi yêu cầu bạn làm là hoàn thành hai mô-đun trong hướng dẫn Năng lực thông tin (Information Literacy), như nêu trong phần Tài liệu đọc bắt buộc của bài học này. Hãy đảm bảo hoàn thành các mô-đun sau:

\begin{itemize}
\item
  TU2: Xác định các nguồn tài liệu cần thiết
\item
  TU5: Đánh giá các nguồn tài liệu
\end{itemize}

Với khối lượng và sự đa dạng khổng lồ của thông tin trên internet, việc đánh giá các nguồn trực tuyến quan trọng hơn bao giờ hết. Bạn có thể tin thông tin từ một blog không? Còn thông tin được đăng như một câu trả lời cho câu hỏi của ai đó thì sao?

Khi chuẩn bị thu thập dữ liệu và tài liệu nghiên cứu cho các nghiên cứu thực nghiệm của mình, bạn có thể gặp các bài báo và tiêu đề tin tức mà bạn dùng trong phân tích cảm xúc hoặc phân tích văn bản. Hãy lưu ý rằng có rất nhiều tin giả trên internet.

\begin{enumerate}
\def\labelenumi{\arabic{enumi}.}
\setcounter{enumi}{2}
\item
  Các nguồn không đáng tin cậy
\end{enumerate}

Báo chí đã thay đổi mạnh mẽ trong vài thập kỷ qua cùng với sự trỗi dậy của internet, điện thoại thông minh và mạng xã hội. Sự phát triển của các kênh tin tức phát sóng 24 giờ đã khuyến khích các kênh tin tức truyền hình cáp áp dụng những tiêu đề giật gân, đồng thời đề cao quan điểm thiên về đảng phái của chủ sở hữu và cân bằng nhu cầu cũng như yêu cầu của các nhà quảng cáo. Truyền hình và mạng xã hội đã vượt qua báo in để trở thành nguồn tin tức, nên báo chí và báo chí địa phương phải vật lộn để tồn tại. Kết quả là mọi người phải tìm đến các nguồn khác để có tin tức, và những nguồn này thường có thể có xung đột lợi ích cần được xem xét khi đánh giá liệu nội dung đưa tin có đáng tin cậy và uy tín hay không. Báo chí chịu rất nhiều áp lực phải đưa tin trước, bất kể phóng viên đã nắm đủ các sự kiện hay chưa. Do đó, các bản tin ban đầu thường sai, và các kênh truyền hình hoặc ấn phẩm có thể phải đăng đính chính hoặc rút lại tin sau đó.

Đây đều là những thực tế quan trọng cần ghi nhớ khi tìm kiếm nguồn trên internet. Điều này không có nghĩa là mọi kênh tin tức truyền hình cáp đều không đáng tin hay bạn không bao giờ nên tin một blog. Thay vào đó, điều cần rút ra là hãy tự tìm hiểu kỹ. Khi quyết định một nguồn có đáng tin cậy hay không, hãy đặt một vài câu hỏi then chốt:

Ai là người viết văn bản này? Họ có danh tiếng là học giả hoặc chuyên gia đáng tin cậy trong lĩnh vực hoặc chủ đề này không? Đây có phải là một tổ chức đáng tin cậy, đã chứng tỏ cam kết đưa tin chính xác các sự kiện không? Có thiên lệch (bias) chính trị hoặc xung đột lợi ích tiềm ẩn nào mà bạn cần lưu ý không?

-Văn bản này được viết khi nào? Nguồn này có lỗi thời hoặc không còn phù hợp với chủ đề bạn đang nghiên cứu không? Nguồn này phản hồi hoặc tiếp nhận những diễn biến hiện đại liên quan đến chủ đề như thế nào?

-Lập luận, logic và cách suy luận của văn bản có vững vàng không? Thông tin trong văn bản có dễ bị bác bỏ không? Phần lớn các sự kiện hoặc số liệu thống kê có thể được chứng minh bằng một nguồn khác không? Tác giả có trích dẫn các nguồn đáng tin cậy làm bằng chứng không?

-Phong cách hoặc giọng điệu của văn bản là gì? Văn bản mang tính học thuật và khách quan hơn, hay mang tính cảm xúc và kích động? Văn bản có đang cố thuyết phục bạn mua/sử dụng/đọc/xem/làm điều gì đó không?

Khi hoàn thành phần hướng dẫn về năng lực thông tin, bạn sẽ học được các kỹ năng và chiến lược cần thiết để đánh giá một nguồn và trả lời những câu hỏi này.

Giả sử bạn đang xây dựng một chiến lược giao dịch bằng các thuật toán học máy thực hiện phân tích cảm xúc. Các phương pháp này đòi hỏi tiêu đề, bài báo tin tức và thậm chí cả các bài đăng trên mạng xã hội. Sự tồn tại của tin giả có thể biến những mô hình này thành những chiếc hộp "garbage in, garbage out" (đầu vào rác, đầu ra rác). Hoặc có thể, khi thực hiện dự án capstone, bạn quyết định dùng một nguồn mà hóa ra là giả. Điều này có thể gây ra những hệ quả tai hại cho kết quả nghiên cứu của bạn. Trong thời đại thông tin sai lệch, khả năng nhận diện (và tránh) các nguồn không đáng tin cậy quan trọng hơn bao giờ hết.

\begin{enumerate}
\def\labelenumi{\arabic{enumi}.}
\setcounter{enumi}{3}
\item
  Tư duy phản biện về tin giả
\end{enumerate}

Vui lòng đọc bài ``The Use of Critical Thinking to Identify Fake News: A Systematic Literature Review'' của Machete và Turpin. Bài viết tổng quan các nghiên cứu về cách tư duy phản biện có thể được dùng để nhận diện tin giả. Những bài báo thảo luận rõ ràng về tư duy phản biện cho thấy phần lớn sinh viên không có đủ kỹ năng tư duy phản biện để nhận diện tin tức giả.

Sau khi đọc xong bài viết, hãy áp dụng những gì bạn đã học từ khóa học về năng lực thông tin để xác định một nội dung có phải là tin giả hay không.

\begin{enumerate}
\def\labelenumi{\arabic{enumi}.}
\item
  Thông tin có chính xác không, xét theo việc đối chiếu với các nguồn khác?
\item
  Bài viết có do một người có thẩm quyền viết không? Người đó có được giới đồng nghiệp công nhận không?
\item
  Thông tin có đầy đủ không? Cần lưu ý những câu chuyện bỏ sót thông tin quan trọng.
\item
  Câu chuyện có kịp thời không? Có bản cập nhật hoặc tin tức nào khác cho thấy thông tin có thể đã lỗi thời không?
\item
  Thông tin được trình bày như tiêu đề câu khách (clickbait), chỉ nhằm thu hút người đọc, hay là một bản tin thực sự?
\end{enumerate}

Giả sử trong quá trình thẩm định, bạn trả lời "không" cho một số câu hỏi trên. Khi đó, câu hỏi cần đặt ra là "câu chuyện này có được lặp lại ở nơi khác không, và đã được các nguồn đáng tin cậy xác nhận hoặc kiểm chứng chưa?" Bất kỳ nghi ngờ nào của bạn đều có thể khiến bài viết đó không trở thành một phần dữ liệu cho phân tích cảm xúc của bạn.

\begin{enumerate}
\def\labelenumi{\arabic{enumi}.}
\setcounter{enumi}{4}
\item
  Kết luận
\end{enumerate}

Rõ ràng, việc tạo ra tin giả là vi phạm thực hành đạo đức tốt. Là một FE, bạn có thể nâng cao năng lực thông tin bằng cách học cách đánh giá chất lượng, thẩm quyền, tính toàn diện và tính kịp thời của thông tin bạn tìm thấy trên web. Thông tin này có thể trở thành dữ liệu mà bạn dùng để xây dựng và huấn luyện các mô hình cho FE. Xây dựng năng lực thông tin thực ra cũng sẽ giúp thúc đẩy thực hành đạo đức tốt.

Khi kết thúc học phần này, hãy nhớ trung thực và cảnh giác để tránh tạo ra, lan truyền hoặc đưa tin giả vào các nghiên cứu của bạn. Ngay cả khi bạn dùng AI để xác định đâu là tin giả, hãy chắc chắn rằng bạn xác nhận lại những gì được sử dụng. Câu nói "rác vào, rác ra" (garbage in, garbage out) vẫn đúng, nhưng tình huống có thể tệ hơn nếu chúng ta cho rằng các câu chuyện giả là thật và đưa ra những quyết định quan trọng dựa trên thông tin sai.
